\section{Diffraction about a smooth macroimage}

We use the single-plane perturbative framework of
Oguri and Takahashi~\cite{oguri2022}.  The stochastic halo signal studied
by Choi et al.~\cite{choi2026} motivates the application, but none of the
normalizations below is calibrated to their line-of-sight CDM calculation.
All coordinates in this section are physical lens-plane coordinates.
For angular-diameter distances $D_l,D_s,D_{ls}$ and observed frequency $f$,
define
\begin{equation}
 r_F^2(f)=\frac{cD_lD_{ls}}{2\pi f(1+z_l)D_s},\qquad
 A=\begin{pmatrix}\lambda_1&0\\0&\lambda_2\end{pmatrix},\qquad
 |\mu|=\frac{1}{|\lambda_1\lambda_2|}.
\end{equation}
The eigenvalues may have either sign but must be nonzero.  Locally, the
smooth Fermat potential is quadratic.  Dividing its perturbed Fresnel
integral by the unperturbed one gives
\begin{equation}
 \frac{F}{F_0}=
 \frac{\displaystyle\int d^2x\,
  \exp\left[\frac{i\boldsymbol{x}^{T}A\boldsymbol{x}}{2r_F^2}
                -\frac{i\delta\psi(\boldsymbol{x})}{r_F^2}\right]}
 {\displaystyle\int d^2x\,
  \exp\left[\frac{i\boldsymbol{x}^{T}A\boldsymbol{x}}{2r_F^2}\right]}.
 \label{eq:macro-fresnel}
\end{equation}
The smooth factor has modulus $|F_0|=\sqrt{|\mu|}$ and the appropriate
Morse phase.  The perturbing potential obeys
$\nabla^2\delta\psi=2\delta\kappa$.
We remove the arrival-time shift $-\delta\psi(0)/r_F^2$ from the phase
and retain first order in $\delta\psi$.  With Fourier convention
$\delta\kappa(\boldsymbol{x})=
\int d^2k\,\widetilde\kappa(\boldsymbol{k})
e^{i\boldsymbol{k}\cdot\boldsymbol{x}}/(2\pi)^2$,
the Gaussian integral yields
\begin{align}
 \eta(f)&\equiv e^{i\delta\psi(0)/r_F^2}\frac{F(f)}{F_0(f)}-1
 =\int\frac{d^2k}{(2\pi)^2}\,
 \widetilde\kappa(\boldsymbol{k}) H_A(f,\boldsymbol{k}),\\
 H_A(f,\boldsymbol{k})&=
 \frac{2i}{r_F^2 k^2}
 \left\{\exp\left[-\frac{i r_F^2}{2}
 \boldsymbol{k}^{T}A^{-1}\boldsymbol{k}\right]-1\right\}.
 \label{eq:macro-kernel}
\end{align}
Equation~\eqref{eq:macro-kernel} is Eq.~(21) of
Ref.~\cite{oguri2022}, with their directional magnifications equal to
$\lambda_1^{-1},\lambda_2^{-1}$.  It applies to saddles as well as minima.
For $A=I$, $H=(1-e^{-is})/(is)$, where $s=r_F^2 k^2/2$.

The numerical implementation uses
$q=k r_F(f_{\rm ref})$, $f_{\rm ref}=100\,\mathrm{Hz}$, and
\begin{equation}
 b(\theta)=\frac{\cos^2\theta}{\lambda_1}
           +\frac{\sin^2\theta}{\lambda_2},\qquad
 s=\frac{q^2b(\theta)f_{\rm ref}}{2f},\qquad
 H_A=b(\theta)\,\frac{\sin(s/2)}{s/2}e^{-is/2}.
\end{equation}
This form is continuous at $q=0$ and $b=0$.  At fixed $\boldsymbol{k}$,
$H_A\to b(\theta)$ as $f\to\infty$: a real, frequency-independent
geometric magnification, not a measurable diffraction residual by itself.

\subsection{Projection and real covariance}

For normalized positive frequency weights $w_a$, the relevant inner
product is $\operatorname{Re}\sum_a w_a u_a^*v_a$.
The nuisance directions for an image or an aligned image ratio are
$1,i,if$, corresponding to amplitude, phase, and arrival time.  We project
the real part of $H$ off a constant and its imaginary part off a constant
plus a linear function of $f$:
\begin{align}
 \operatorname{Re}(QH)&=\operatorname{Re}H-
                  \langle\operatorname{Re}H\rangle_w,\\
 \operatorname{Im}(QH)&=\operatorname{Im}H-
                  \langle\operatorname{Im}H\rangle_w
 -(f-\langle f\rangle_w)
  \frac{\langle(f-\langle f\rangle_w)\operatorname{Im}H\rangle_w}
       {\langle(f-\langle f\rangle_w)^2\rangle_w}.
 \label{eq:real-projection}
\end{align}
Complex regression on $\{1,f\}$ would also discard a real amplitude
slope and would therefore remove a physical mode unnecessarily.

Define the homogeneous convergence spectrum by
$\langle\widetilde\kappa(\boldsymbol{k})
\widetilde\kappa^*(\boldsymbol{k}')\rangle
=(2\pi)^2\delta^{(2)}(\boldsymbol{k}-\boldsymbol{k}')P_\kappa(k)$
and $\Delta_\kappa^2(k)=k^2P_\kappa(k)/(2\pi)$.
Writing
$\boldsymbol{v}(k,\theta)=[\operatorname{Re}(QH),
\operatorname{Im}(QH)]^T$, the unweighted joint covariance is
\begin{equation}
 C=\int d\ln k\,\Delta_\kappa^2(k)
       \left\langle\boldsymbol{v}\boldsymbol{v}^{T}\right\rangle_\theta.
 \label{eq:real-covariance}
\end{equation}
The mean-square projected fractional residual is
$\sigma_\perp^2=\operatorname{tr}(WC)$ with
$W=\operatorname{diag}(w_1,\ldots,w_N,w_1,\ldots,w_N)$.
The real field makes the complex fluctuations noncircular; the joint
real covariance retains information lost by keeping only
$\langle\eta\eta^*\rangle$.

For two images, the cross-covariance before projection includes
$H_{A_1}(f,\boldsymbol{k})H_{A_2}^*(g,\boldsymbol{k})
e^{i\boldsymbol{k}\cdot(\boldsymbol{x}_1-\boldsymbol{x}_2)}$.
Their differential covariance is $C_1+C_2-C_{12}-C_{21}$.
Using $C_1+C_2$ therefore requires independent image perturbations;
it is not automatic for all line-of-sight structures.

\subsection{An analytic power-law check}

Suppose $P_\kappa(k)=C_0 k^n$ on all relevant scales.  Consider a fixed
fractional frequency band $f=f_0 u$ with identical weights as a function
of $u$.  Changing radial variable to
$p=k r_F(f_0)\sqrt{|b(\theta)|}$ in
Eq.~\eqref{eq:real-covariance} gives
\begin{equation}
 \sigma_\perp^2(f_0,A)\ \propto\
 f_0^{(n+2)/2}
 \left\langle |b(\theta)|^{(2-n)/2}\right\rangle_\theta,
 \qquad -6<n<2.
 \label{eq:powerlaw-scaling}
\end{equation}
The convergence range follows from $QH=O(k^2)$ at small $k$ and
$QH=O(k^{-2})$ at large $k$.  Thus the RMS scales as
$f_0^{(n+2)/4}$.  Frequency independence occurs only for $n=-2$;
it is not a general prediction for CDM or for two detectors with
different waveform and noise weights.

For $n=-2$, the RMS gain relative to $A=I$ has the closed form
\begin{equation}
 G^2\equiv\frac{\sigma_\perp^2(A)}{\sigma_\perp^2(I)}
 =\left\langle b^2\right\rangle_\theta
 =\frac{3}{8}\left(\lambda_1^{-2}+\lambda_2^{-2}\right)
   +\frac{1}{4\lambda_1\lambda_2}.
 \label{eq:flat-spectrum-gain}
\end{equation}
For $\lambda_2=1$ and $\lambda_1=\pm1/|\mu|$,
$G\simeq\sqrt{3/8}\,|\mu|$ at large $|\mu|$.
This fractional enhancement is separate from the smooth-waveform SNR
gain $\sqrt{|\mu|}$.  The radial and angular numerical integrals are
tested against Eq.~\eqref{eq:flat-spectrum-gain}; finite spectral cutoffs
and changes of slope require the full integral.

\paragraph{Scope and validity.}
These results describe a weak screen at the main lens plane and a
locally quadratic, nonsingular macroimage.  The elliptical Fresnel
widths $r_F/\sqrt{|\lambda_i|}$ must remain within that local region,
and perturbations must remain small enough for the Born expansion.
One cannot apply the same $A^{-1}$ to every foreground and background
screen: their coupling to the macro lens needs a multiplane propagation
calculation.  The calculation also does not follow newly created
microimages or nonlinear caustics.  Stellar fluctuations can dominate
CDM fluctuations in the ground-based band even when both are enhanced
by strong lensing~\cite{oguri2022}.
